% !TeX root = ../main.tex
\subsection{BARCS estimates multivariable effects from guide counts}
\label{sec:overview}

BARCS takes a guide-count matrix, library totals, and a sample design matrix.
For guide $g$ in library $i$, the default model is
\begin{equation}
K_{gi}\mid N_i\sim
\operatorname{BetaBinomial}(N_i,\mu_{gi},\rho_g),
\qquad
\logit(\mu_{gi})=\mathbf{x}_i^\top\boldsymbol{\beta}_g,
\label{eq:barcs-model}
\end{equation}
where $N_i$ is the full mapped-guide total before filtering, $\mu_{gi}$ is the
expected guide proportion, and $\rho_g$ represents overdispersion.
The design vector $\mathbf{x}_i$ can contain continuous predictors, group or
donor indicators, and interactions.  A prespecified coefficient or contrast
therefore addresses a question within one joint fit.  BARCS retains the
sampling principle of CB\textsuperscript{2} \citep{jeong2019beta}, with its own
estimation and finite-sample testing procedure.

The workflow fits guide coefficients and dispersion, computes a $t$ statistic
using residual sample degrees of freedom, and summarizes guide evidence when
the reported unit is a gene.  Gene summaries use a directional Stouffer
statistic and the median guide coefficient.  In the simulation ablations,
BARCS-ST uses unmoderated guide statistics and BARCS-MOD applies empirical-Bayes
variance moderation before the same summary.  Real-data analyses use
unmoderated BARCS.  Control scaling adjusts significance without changing
coefficient estimates; its placement is specified for each benchmark.
The Cas13 comparison calibrates gene summaries using aggregation-matched
control pseudo-genes, whereas the IL2RA diagnostic cross-fits a guide-level
scale.  The null simulations below test whether guide-level calibration
extends to gene summaries.

Full-library totals estimate changes in library share.  Control-based
normalization instead expresses abundance relative to a prespecified control
class, as examined in the interaction simulation.  This deliberate change of
reference is distinct from inadvertently recomputing totals after guide
filtering.  The following analyses evaluate the additional design information,
the effects of these optional steps, and the assumptions needed to interpret
the resulting probabilities.
